\input{../tex-inputs/latex-leseansicht-vorspann}

\section[Rainer Maria Rilke an Arthur Schnitzler, {[}Ende 1901{]}]{L04541 Rainer Maria Rilke an Arthur Schnitzler, {[}Ende 1901{]}}
\nopagebreak\mylabel{L04541v}
\rehead{ }\normalsize\beginnumbering\briefempfaengerindex{Schnitzler, Arthur@\textsc{Schnitzler, Arthur}!zzzRilke, Rainer Maria@\emph{von Rainer Maria Rilke}!1901-12-312@{{[}31. 12. 1901{]}}|(be}
\toendnotes[C]{\smallbreak\pagebreak[2]}
\correspDesc{Versand  durch Rainer Maria Rilke am [31. 12. 1901] in ?? [Bauernhaus von Clara und Rainer Maria Rilke in Westerwede]
\newline{}Erhalt  durch Arthur Schnitzler am [7. 1. 1902?] in Frankgasse 1}\toendnotes[C]{\smallbreak}
\Standort{DLA, A:Rilke-Archiv Gernsbach/Kopien und Abschriften, HS.2022.81.}
\physDesc{Bildpostkarte, Fotokopie, 271 Zeichen
\newline{}Handschrift: schwarze Tinte, deutsche Kurrent
\newline{}Schnitzler: mit Bleistift datiert: »Ende 901.« 
\newline{}Ordnung: 1) Die Fotokopien wurden zwischen 30. 12. 1930 und
                                 14. 1. 1931 durch Carl Sieber\pwindex{Sieber, Carl 1897 – 1945@\textsc{Sieber, Carl} (1897 – 1945)|pw} veranlasst, bevor die
                                 Originale an Schnitzler
                                 zurückgestellt wurden. Die Adressseite der Karte wurde nicht abgelichtet.  2) Das Original dieser Karte ist nicht erhalten. Es dürfte von Heinrich Schnitzler\pwindex{Schnitzler, Heinrich 9.\,8.\,1902 Hinterbrühl – 12.\,7.\,1982 Wien@\textsc{Schnitzler, Heinrich} (9.\,8.\,1902 Hinterbrühl – 12.\,7.\,1982 Wien)|pw} am 15. 8. 1936 an
                                  Paul Marx\pwindex{Marx, Paul 21.\,7.\,1879 Wien – 30.\,10.\,1956 ebd.@\textsc{Marx, Paul} (21.\,7.\,1879 Wien – 30.\,10.\,1956 ebd.)|pw} geschenkt worden sein.}
\buchAbdrucke{\weitereDrucke{\emph{Rainer Maria Rilke und Arthur Schnitzler. Ihr Briefwechsel.}. Mit Anmerkungen herausgegeben von Heinrich Schnitzler In: \emph{Wort und Wahrheit}, Jg. 13, Nr. 4, April 1958, S. 282–298, hier S. 289.} }\toendnotes[C]{\smallbreak}
\pstart
           \centering{}{\pb}\textcolor{gray}{\textbf{Torfſtich im Moor.}}\pend
           
\pstart
           \centering{}\textcolor{gray}{\textbf{Gruss sendet}}\pend
           \vspace{1em}
\pstart
           \noindent{}{\pb}Grüße zum kommenden Jahr! Von mir und meiner Frau, \textsc{Clara Westhoff\pwindex{Westhoff, Clara 21.\,11.\,1878 Bremen – 9.\,3.\,1954 Fischerhude@\textsc{Westhoff, Clara} (21.\,11.\,1878 Bremen – 9.\,3.\,1954 Fischerhude)|pw}}! Ich werde leider \label{K_L04541-1v}\edtext{zum 4.\eventindex{Deutsches Theater Berlin@\textbf{Deutsches Theater Berlin}!Uraufführung von Lebendige Stunden, 4.1.1902@Uraufführung von Lebendige Stunden, 4.1.1902|pwv} nicht in \textsc{Berlin\oindex{Berlin@\textbf{Berlin}|pw}}}{\lemma{\textnormal{\emph{zum 4. nicht in Berlin}}}\Cendnote{\textnormal{Vgl. Uraufführung von Lebendige Stunden, 4.1.1902. In: A. S.: \emph{Kulturveranstaltungen}, \url{https://schnitzler-kultur.acdh.oeaw.ac.at/pmb35104.html}.}}}\label{K_L04541-1}{ }ſein können. Nicht{ }ſelbſt: aber mit Gedanken und
               Wünſchen! Dank für Ihren lieben \label{K_L04541-2v}\edtext{Brief}{\lemma{\textnormal{\emph{Brief}}}\Cendnote{\textnormal{XXXX Auszeichnungsfehler: Dokument L04344 nicht gefunden.}}}\label{K_L04541-2}, auf den ich noch ausführlich zurückkomme im neuen
               Jahre:\pend
           
\pstart
           der Ihre:{\\[\baselineskip]}\spacefill\mbox{Rainer Maria Rilke}\pend
           \leftskip=0em{}\selectlanguage{ngerman}\endnumbering\briefempfaengerindex{Schnitzler, Arthur@\textsc{Schnitzler, Arthur}!zzzRilke, Rainer Maria@\emph{von Rainer Maria Rilke}!1901-12-312@{{[}31. 12. 1901{]}}|)be}\mylabel{L04541h}
\begin{anhang}
\end{anhang}\newcommand{\dateiname}{L04541}\newcommand{\titel}{Rainer Maria Rilke an Arthur Schnitzler, [Ende 1901]}\newcommand{\editorInnen}{Herausgegeben von Jahnke, SelmaMüller, Martin Anton}\input{../tex-inputs/latex-leseansicht-abspann}
